\scanpage{51}
arête sss $a$ n'est contenu dans aucun
des $S_i$. D'autre part, si $u$\inserted{$\subset S$} est un
$n$-uple, alors
\[
u\longrightarrow I
\]
est injectif \inserted{(i.e. $u\cap E_i$ de card $\leq1$ pour tt $i$).} Dire que c'est aussi surjectif
(ce qui est vérifié dans le cas d'une $n$-alité)
signifie ceci

(iii) $\forall s\in S,\ x\in R$, existe $s'\in S$ tel que
$\{s,s'\}$ ne soit pas une arête (p.ex. $s'=s$) et $(s',x)\in R$
(NB en vertu de (ii) cet $s'$ est unique).
% slanted
\margin{implique que tous les $R(x)$ sont de $=$ cardinal $n=\operatorname{card}I$}

Alors on a
\[
\varphi_i:E\longrightarrow\deleted{\unsure{S_i}}E_i
\]
en associant à tout $x\in E$ \deleted{le système}
\deleted{des} $\varphi_i(x)\in E_i$ tel que $\{\varphi_i(x)\}=S(x)\cap E_i$

(iv) \deleted{\unsure{$\exists s,s'$} \illegible{1}} $\exists x\in E$ tel que $R(x)$ de card $\geq2$.
\[
(\Longleftrightarrow\operatorname{card}I\geq2)
\]

\uline{NB} Dans le cas contraire, tous les $R(x)$
sont de cardinal $1$ (On suppose $S\ne\unsure{0|\varnothing}$)
et $R$ est le graphe d'une
application $E\mathrel{\unsure{-|\to}}S$

\begin{center}\rule{1.2cm}{.6pt}\end{center}
% archive sheet number: 50
